\clearpage
\section{Заключение}
\label{sec:conclusion}

\subsection{Выполненные задачи}

\begin{enumerate}
\item Рассмотрены подходы к определению результата молодой компании и разграничены прогнозирование, описание зависимостей и причинная интерпретация.
\item Сформирована исходная группа активных компаний, разделены ранние признаки и поздние статусы. Происхождение файлов описано в разделе~\ref{sec:data}, состав когорты и результаты сопоставления --- в разделе~\ref{sec:eda}. Для позднего файла документировано основание отнесения к августовской версии 2026 года.
\item Выполнен разведочный анализ полноты, состава и распределений признаков; отдельно исследованы неизвестные исходы и группа IK12.
\item Проведена вложенная проверка постоянного прогноза, логистических моделей и CatBoost. В каждой из двух схем выполнено пять повторов; разнообразие групповых разбиений проверено отдельно.
\item Оценены добавление и исключение блоков признаков, сравнение полной модели с моделью с тегами, разброс между наборами и условные границы при повторных выборках компаний и наборов.
\item Согласованы метрики и PR-графики, проверены веса классов при одинаковом критерии настройки, интерпретированы коэффициенты частых тегов. Проведено повторное обучение без IK12 и без публичных компаний.
\end{enumerate}

\subsection{Основные результаты}

Исходная группа включает 3\,110 компаний, активных 13 июля 2023 года. В августовской версии 2026 года известны статусы 3\,046 компаний; 272 имеют положительную метку, что составляет 8,93\%. Из них 265 поглощены, поэтому целевая метка преимущественно характеризует поглощение и не является общей мерой экономической успешности.

Из 64 неизвестных исходов 43 относятся к IK12. Поэтому почти полное сопоставление остальных компаний не обеспечивает представительности для этого нестандартного набора. Границы 8,75--10,80\% характеризуют возможную общую долю положительных исходов, но не диапазон качества моделей.

Описательное сравнение показало перекрытие распределений размера команды: медианы равны десяти сотрудникам у сохранивших активность, пяти --- у ставших неактивными и восьми --- у поглощённых компаний. Корреляция Спирмена между возрастом и размером команды составляет 0,58. Эти результаты описывают выборку, не устанавливая причинных эффектов или дополнительного прогнозного вклада.

LogLoss постоянного прогноза составляет \resConstantLoss{}, полной логистической модели --- \resFullLoss{}; относительное снижение равно \resRelativeGain{}\%. ROC-AUC полной модели равна \resFullAUC{}, AP --- \resFullAP{}. Результат характеризует небольшой прогностический сигнал в рассматриваемых открытых сведениях.

Добавление тематических меток изменило LogLoss относительно модели с командой на 0,00510 в сторону уменьшения. Условный диапазон повторной выборки наборов целиком находится выше нуля. Это оценка блока признаков в выбранной процедуре обучения, а не доказательство причинного действия тегов.

Преимущество полной модели перед моделью с тегами не установлено устойчиво: условный диапазон по наборам включает ноль.

Среднее улучшение от добавления тегов сохраняется при исключении IK12 и публичных компаний. Проверки чувствительности используют ту же исходную базу и не являются внешней валидацией.

При одинаковом подборе по AP взвешенная модель имеет среднюю AP выше, а точность верхней десятой части выше, чем вариант без весов. Эти наблюдаемые различия не трактуются как доказательство универсального эффекта взвешивания.


В верхней десятой части списка модели с тегами положительную метку имеют \resTagsPrecisionPercent{}\% компаний, Lift составляет \resTagsLift{}. Такой результат описывает обогащение отобранной группы редким исходом; он не является оценкой доходности или готовой рекомендацией по инвестированию. Устойчивость знаков коэффициентов частых тегов относится к выбранной регуляризованной модели и не заменяет отдельной проверки полезности каждой метки.

\subsection{Ограничения исследования}

\begin{itemize}
\item Выборка ограничена компаниями YC, представленными в сторонней выгрузке и активными в июле 2023 года. Результаты не описывают все компании акселератора или весь рынок.
\item Поздний файл скачан автором в середине августа 2026 года. Дата 01.08.2026 обозначает предполагаемую августовскую версию по графику обновления источника; номер версии и точные моменты сбора строк не сохранены. Статусы могут обновляться с задержкой.
\item Два среза не раскрывают даты событий, промежуточные состояния, условия сделок или доходность. Интервал сопоставления общий для компаний; различие возраста не означает различной длительности последующего наблюдения.
\item Неизвестные исходы распределены неравномерно. Исключение IK12 в дополнительном анализе не восстанавливает статусы отсутствующих компаний.
\item Признаки измерены в июле 2023 года, а не при основании или поступлении в YC. Теги отражают как деятельность, так и правила ведения каталога.
\item Внешние части относятся к вложенной проверке внутри данной совокупности. Повторы используют одни и те же компании; условные диапазоны не учитывают переобучение и весь процесс выбора моделей.
\item Набор моделей и проверок уточнялся после изучения данных. Сравнения нескольких блоков и интерпретации отдельных тегов являются разведочными; формальные подтверждающие тесты не проводились.
\end{itemize}

\subsection{Выводы и рекомендации}

Исследование оценивает полезность блоков доступных сведений для конкретного наблюдаемого исхода. Необнаруженное улучшение в отдельной конфигурации не доказывает бесполезности признака во всех моделях или для других исходов. При сравнении сложной и простой модели следует учитывать размер и устойчивость выигрыша, а не только минимальную среднюю ошибку.

Для предварительного ранжирования необходимы независимая проверка качества и калибровки, определение стоимости проверки компании и последствий ошибок. Само по себе преобладание отрицательных исходов среди отобранных компаний при редком событии не исключает практической полезности. Доказательство причинности не обязательно для прогнозного ранжирования, но необходимо для рекомендаций предпринимателю изменить состав команды или позиционирование ради улучшения результата.

Следующий эксперимент следует зафиксировать до наблюдения новых исходов: определить дату формирования признаков, горизонт, набор компаний, модели, метрики и правила отбора. Его нужно проверить на новом календарном периоде, а затем отдельно на компаниях вне YC. После изучения всей текущей совокупности выделение части этих же данных нельзя представлять как полностью независимую подтверждающую проверку.

Данные с датами поглощений, выхода на биржу, прекращения деятельности и финансирования позволят исследовать время до события и конкурирующие исходы. Для текущей постановки с общим поздним срезом бинарная классификация остаётся допустимой; её ограничения связаны прежде всего со смыслом статуса, качеством регистрации и переносимостью результата.
\clearpage
